% !PRINTED 485
\weberpage{485}
%<<CONT
Es handelt sich nämlich noch um den Nachweis, dass man über $\eta_0, \eta_0', \eta_0''$ (§.\,130) so verfügen kann, dass die Functionen $X_1, X_2, X_3$ nicht identisch verschwinden. Dazu müssen wir die Bildungsweise der Grössen $X$ etwas genauer betrachten.

Setzen wir in (9) zufolge (8) (§.\,130):
\[
\frac{\partial^3 f}{\partial \xi_1^3} = 6\,\xi_3, \quad \frac{\partial^3 f}{\partial \xi_1^2 \xi_2} = 0, \quad \frac{\partial^3 f}{\partial \xi_1^2 \partial \xi_3} = 6\,\xi_1, \quad \frac{\partial^3 f}{\partial \xi_1 \partial \xi_2 \partial \xi_3} = 0, \dots, % [sic] source omits the second \partial before \xi_2
\]
so ergiebt sich

\weq{1}{\begin{aligned}
X_1 &= A_{2,2,2} + 3\,A_{1,1,3}, \quad X_2 = A_{3,3,3} + 3\,A_{2,2,1}, \\
X_3 &= A_{1,1,1} + 3\,A_{3,3,2}.
\end{aligned}}

Nun ist aber ferner nach (4) und (6) (§.\,130):

\weq{2}{\Psi =
\begin{vmatrix}
p_1 x_1 + q_3 x_2 + q_2 x_3, & q_3 x_1 + p_2 x_2 + q_1 x_3, & q_2 x_1 + q_1 x_2 + p_3 x_3 \\
p_1' x_1 + q_3' x_2 + q_2' x_3, & q_3' x_1 + p_2' x_2 + q_1' x_3, & q_2' x_1 + q_1' x_2 + p_3' x_3 \\
p_1'' x_1 + q_3'' x_2 + q_2'' x_3, & q_3'' x_1 + p_2'' x_2 + q_1'' x_3, & q_2'' x_1 + q_1'' x_2 + p_3'' x_3
\end{vmatrix}.}

Dies lässt sich leicht nach Potenzen und Producten der $x$ ordnen, und wenn wir also die Bezeichnung gebrauchen
\[
(p_1, q_3, q_2) =
\begin{vmatrix}
p_1, & q_3, & q_2 \\
p_1', & q_3', & q_2' \\
p_1'', & q_3'', & q_2''
\end{vmatrix} \cdots,
\]
so erhält man
\[
\begin{aligned}
A_{1,1,1} &= (p_1, q_3, q_2), & 3\,A_{3,3,2} &= (q_3, q_1, p_3) + (q_2, p_2, p_3) \\
A_{2,2,2} &= (q_3, p_2, q_1), & 3\,A_{1,1,3} &= (p_1, q_1, q_2) + (p_1, q_3, p_3) \\
A_{3,3,3} &= (q_2, q_1, p_3), & 3\,A_{2,2,1} &= (q_3, p_2, q_2) + (p_1, p_2, q_1),
\end{aligned}
\]
und daraus nach (1)

\weq{3}{\begin{aligned}
X_1 &= (q_3, p_2, q_1) + (p_1, q_1, q_2) + (p_1, q_3, p_3) \\
X_2 &= (q_2, q_1, p_3) + (q_3, p_2, q_2) + (p_1, p_2, q_1) \\
X_3 &= (p_1, q_3, q_2) + (q_3, q_1, p_3) + (q_2, p_2, p_3).
\end{aligned}}

Diese Functionen $X_1, X_2, X_3$ sind, wie aus dem oben Bemerkten folgt [§.\,130, (5)], trilineare Formen der drei Variablenreihen $\eta_r, \eta_r', \eta_r''$, deren Co\"efficienten rational durch die siebente Einheitswurzel $\varepsilon$ ausgedrückt werden können. Die Co\"efficienten dieser Formen sind aber gewiss nicht alle gleich Null. Denn nach (3) kann man die Grössen $p, q$ so annehmen, dass $X_1, X_2, X_3$ nicht gleich Null werden, und die 21 Grössen $\eta_r, \eta_r', \eta_r''$ lassen sich aus §.\,130, (5) so bestimmen, dass die 18 Grössen $p, q$ (und ausserdem die drei Summen $\Sigma \eta_r$) beliebig vorgeschriebene Werthe
%CONT>>
